% 大熊座（综述）
% license CCBYSA3
% type Wiki

本文根据 CC-BY-SA 协议转载翻译自维基百科\href{https://en.wikipedia.org/wiki/Ursa_Major}{相关文章}。

大熊座（Ursa Major），亦称“巨熊座”，是位于北天的一座星座，其相关神话很可能可以追溯至史前时代。其拉丁名称意为“较大的（或更大的）熊”，用以指代并区别于附近的小熊座（Ursa Minor）。$^{[2]}$ 在古代，大熊座是公元2世纪托勒密所列出的最初48个星座之一，其编制基础源自更早的希腊、埃及、巴比伦与亚述天文学家的研究成果。$^{[3]}$在现代被正式划定的88个星座中，大熊座按面积排名第三。

大熊座最为人熟知的是由其七颗主星构成的星群，这一星群在不同文化中被称为“大斗”“马车”“查理的车”或“犁”等名称。尤其值得注意的是，“大斗”的恒星排列形态在视觉上与“小斗”相呼应。其中的两颗恒星——天枢与天璇（分别为大熊座$\alpha$星与$\beta$星，Dubhe 与 Merak）——可作为导航指示星，用来指向当前的北极星，即位于小熊座中的北极星（Polaris）。

大熊座及其所包含或交叠的星群在世界诸多文化中具有重要意义，常被视为“北方”的象征。阿拉斯加州旗帜上对大熊座的描绘，便是这一象征在现代的一个实例。

在北半球的大多数地区，大熊座全年可见，并在中北纬度地区呈现为拱极星座。从南半球的温带纬度看，其主要星群不可见，但星座的南部区域仍可被观测到。
\subsection{特征}  
大熊座覆盖约1279.66平方度，约占全天空的3.10\%，是第三大星座。$^{[4]}$ 1930年，欧仁·德尔波特为其划定了国际天文学联合会（IAU）认可的正式星座边界，将其定义为一个由28条边组成的不规则多边形。在赤道坐标系中，该星座的赤经范围介于$08^{\mathrm h}08.3^{\mathrm m}$至$14^{\mathrm h}29.0^{\mathrm m}$之间，赤纬范围介于$+28.30^\circ$至$+73.14^\circ$ 之间。$^{[5]}$大熊座与八个星座相邻：北方与东北方为天龙座，东方为牧夫座，东与东南为猎犬座，东南为后发座，南方为狮子座与小狮座，西南为天猫座，西北为鹿豹座。国际天文学联合会于1922年采用其三字母缩写“UMa”。$^{[5]}$
\subsection{特征}  
\subsubsection{星群}
\begin{figure}[ht]
\centering
\includegraphics[width=6cm]{./figures/f2d88b83474133e7.png}
\caption{大熊座与北极星，以及北斗七星中明亮恒星的名称} \label{fig_UrMa_1}
\end{figure}
\begin{figure}[ht]
\centering
\includegraphics[width=6cm]{./figures/3eb25f5a42cf8f25.png}
\caption{肉眼所见的大熊座星座} \label{fig_UrMa_2}
\end{figure}
大熊座中七颗明亮恒星的轮廓构成了一个星群，在美国和加拿大被称为“北斗”（Big Dipper），而在英国则被称为“犁”（Plough）$^{[6]}$，或在历史上称为“查理之车”（Charles' Wain）$^{[7]}$。这七颗恒星中有六颗的视星等为二等或更亮，使其成为夜空中最为人熟知的恒星图案之一。$^{[8][9]}$正如这些常见名称所暗示的那样，其形状常被认为类似于一把勺子、一种农用犁具，或一辆马车。在大熊座的语境中，它们通常被绘制为巨熊的后躯与尾巴。从“勺斗”部分开始，沿顺时针方向（即天空中的东方方向）延伸至柄部，这些恒星依次为：
\begin{itemize}
\item 天枢（Dubhe，“熊”），视星等为1.79，是全天第35亮的恒星，也是大熊座中第二亮的恒星。  
\item 天璇（Merak，“熊的腰部”），视星等为2.37。  
\item 天玑（Phecda，“大腿”），视星等为2.44。  
\item 天权（Megrez），其名称意为“尾巴的根部”，指其位于熊的躯体与尾巴的交汇处（或北斗中斗身与柄部的连接点）。  
\item 玉衡（Alioth），其名称并非指熊，而是意为“黑马”，这是由原名讹变并误用于与其名称相近的伴星开阳（Alcor）所致，后者是肉眼可见的双星伴侣，围绕着开阳主星运行。$^{[10]}$ 玉衡是大熊座中最亮的恒星，也是全天第33亮的恒星，视星等为1.76。它同时也是化学丰度异常的 Ap 型恒星中最明亮的一颗，这类磁性恒星的某些化学元素会出现异常贫化或富集，并且随着恒星自转而表现出变化。$^{[10]}$
\item 开阳（Mizar，大熊座$\zeta$星），位于北斗柄端向内数第二颗，是大熊座中第四亮的恒星。“Mizar”一名意为“腰带”。它与其视伴星辅星（Alcor，大熊座80号星）共同构成著名的双星系统，在阿拉伯传统中被称为“马与骑士”。  
\item 摇光（Alkaid，大熊座$\eta$星），位于熊尾的末端。其视星等为1.85，是大熊座中第三亮的恒星。$^{[11][12]}$
\end{itemize}
除天枢与摇光之外，北斗七星中的其他恒星都具有指向人马座方向的共同自行运动。人们还识别出少数具有相同运动特征的其他恒星，它们共同被称为“大熊座移动星群”（Ursa Major Moving Group）。
\begin{figure}[ht]
\centering
\includegraphics[width=6cm]{./figures/58bc21206ce95031.png}
\caption{大熊座与小熊座相对于北极星的位置关系} \label{fig_UrMa_3}
\end{figure}
天璇（Merak，大熊座$\beta$星）与天枢（Dubhe，大熊座$\alpha$星）被称为“指极星”，因为它们有助于寻找北极星（亦称北辰星或极星）。通过目视从天璇指向天枢画一条直线（记为1个单位），并沿同一方向继续延伸5个单位，视线便会落在北极星上，从而准确指示真北方向。

在阿拉伯文化中，还识别出另一组星群，被视为一只跳跃羚羊留下的三对足迹。$^{[13]}$ 这是一列位于星座南部边界的三对恒星。从东南至西南依次为：“第一次跳跃”，由$\nu$与$\xi$大熊座（分别为 Alula Borealis 与 Alula Australis）组成；“第二次跳跃”，由 $\lambda$与$\mu$ 大熊座（Tania Borealis 与 Tania Australis）组成；以及“第三次跳跃”，由 $\iota$与$\kappa$ 大熊座（Talitha 与 Alkaphrah）组成。同一星群在中文中称为“三台”$^{[13]}$，在印度传统中称为 Trivikrama$^{[14]}$，两者均意为“三步”。
\subsubsection{其他恒星}  
W 大熊座是接触双星变星的一类原型，其视星等在 7.75 与 8.48 之间变化。

47 大熊座是一颗类太阳恒星，拥有一个由三颗行星组成的系统。$^{[15]}$其中，47 大熊座 b 于1996年发现，其公转周期为1078天，质量为木星的2.53倍。$^{[16]}$47 大熊座 c 于2001年发现，其公转周期为2391天，质量为木星的0.54倍。$^{[17]}$47 大熊座 d 于2010年发现，其周期尚不确定，介于8907至19097天之间，质量为木星的1.64倍。$^{[18]}$该恒星的视星等为5.0，距离地球约46光年。$^{[15]}$

恒星 TYC~3429-697-1（$9^{\mathrm h}40^{\mathrm m}44^{\mathrm s}$，$48^\circ14'2''$），位于 $\theta$ 大熊座以东、北斗以西南的位置，被认定为美国特拉华州的州星，并被非正式地称为“特拉华钻石”。$^{[19]}$
\subsubsection{深空天体}
\begin{figure}[ht]
\centering
\includegraphics[width=6cm]{./figures/a11c8340dbe9123a.png}
\caption{旋涡星系} \label{fig_UrMa_4}
\end{figure}
大熊座中分布着多座明亮的星系，其中包括位于熊头上方的一对星系——梅西耶81（M81，夜空中最明亮的星系之一）与梅西耶82（M82），以及位于摇光（Alkaid）东北方向的旋涡星系（M101）。该星座中还包含旋涡星系梅西耶108（M108）与梅西耶109（M109）。明亮的行星状星云——猫头鹰星云（M97）则位于北斗斗身的底部附近。

M81 是一座几乎正面对着地球的旋涡星系，距离地球约1180万光年。与大多数旋涡星系类似，它的核心由老年恒星构成，而旋臂中则充满年轻恒星与星云。M81 与 M82 同属距离本星系群最近的一个星系团。

M82 是一座几乎侧向可见的星系，正与 M81 发生引力相互作用。它是夜空中最明亮的红外星系。$^{[20]}$2014年1月21日，人们在 M82 中观测到了一次视星等显著的 Ia 型超新星 SN~2014J。$^{[21]}$

M97，又称猫头鹰星云，是一座距离地球约1630光年的行星状星云，其视星等约为10。它于1781年由皮埃尔·梅尚发现。$^{[22]}$

M101，又称旋涡星系，是一座正面对地球的旋涡星系，距离约2500万光年，由皮埃尔·梅尚于1781年发现。其旋臂中存在大规模恒星形成区域，并伴随强烈的紫外辐射。$^{[20]}$该星系的综合视星等为7.5，可通过双筒望远镜或天文望远镜观测，但肉眼不可见。$^{[23]}$

NGC~2787 是一座距离约2400万光年的透镜状星系。与多数同类星系不同的是，它在中心具有明显的棒状结构，并且拥有一个由球状星团构成的晕，表明其具有较高的年龄与相对稳定性。$^{[20]}$

NGC~2950 是一座透镜状星系，距离地球约6000万光年。

NGC~3000 是一对双星，但在星表中被编入星云型天体。

NGC~3079 是一座位于约5200万光年之外的星暴旋涡星系。其中心存在一个马蹄形结构，暗示着超大质量黑洞的存在，该结构由黑洞产生的超风所形成。$^{[20]}$

NGC~3310 是另一座星暴旋涡星系，距离地球约5000万光年。其明亮的白色外观源于异常旺盛的恒星形成速率，这一过程始于约一亿年前的一次星系并合。对该星系及其他星暴星系的研究表明，其星暴阶段可持续数亿年，远长于以往的假设。$^{[20]}$

NGC~4013 是一座侧向可见的旋涡星系，距离地球约5500万光年。它具有显著的尘埃带，并包含多个清晰可见的恒星形成区域。$^{[20]}$

I~Zwicky~18 是一座年轻的矮星系，距离地球约4500万光年。作为可见宇宙中已知最年轻的星系，I~Zwicky~18 的年龄约为400万年，仅为太阳系年龄的约千分之一。它内部充满恒星形成区，以极高的速率生成大量炽热、年轻的蓝色恒星。$^{[20]}$

哈勃深场（Hubble Deep Field）位于大熊座$\delta$星的东北方向。
\subsubsection{流星雨}  
$\alpha$大熊座流星雨（Alpha Ursae Majorids）是该星座中的一场较小型流星雨。$^{[24]}$它们可能由彗星 C/1992~W1（Ohshita）所引发。$^{[24][25]}$

$\kappa$ 大熊座流星雨（Kappa Ursae Majorids）是一场新近发现的流星雨，其极大期位于11月1日至11月10日之间。$^{[26]}$

10月大熊座流星雨（October Ursae Majorids）于2006年由日本研究人员发现。它们可能由一颗长周期彗星所致。$^{[27]}$该流星雨的极大期介于10月12日至10月19日之间。$^{[28]}$
\subsubsection{系外行星}  
HD~80606 是一颗类太阳恒星，位于一个双星系统中，与其伴星 HD~80607 围绕共同的质心运行，两者的平均间距约为1200~AU。2003年的研究表明，其唯一已知行星 HD~80606~b 是一颗未来的“热木星”，模型显示它最初在距离其母星约5~AU、且轨道几乎垂直的轨道上演化而来。该行星质量约为木星的4倍，预计最终会在 Kozai 机制作用下演化为一条更加圆形、且与恒星自转轴更为对齐的轨道。然而在当前阶段，它仍处于一条极端高偏心率的轨道上，其远日点距离约为1个天文单位，而近日点则仅约为6个恒星半径。$^{[29]}$
\subsection{历史}
\begin{figure}[ht]
\centering
\includegraphics[width=6cm]{./figures/e4987f34375d0ea2.png}
\caption{约1700年刻于石上的大熊座，地点为苏格兰法夫郡克雷尔（Crail）} \label{fig_UrMa_5}
\end{figure}
大熊座已被重建为一个印欧语系的星座。$^{[30]}$它是公元2世纪天文学家托勒密在其《天文学大成》（Almagest）中所列出的48个星座之一，托勒密将其称为 Arktos Megale（“大熊”）。$^{[a]}$
\subsection{神话}
大熊座这一星座被多个彼此关联的文明视为一只熊，通常是雌熊。$^{[32][33]}$这可能源于一个共同的口述传统，即“宇宙狩猎”神话，其历史可追溯至一万三千年以上。$^{[34]}$朱利安·迪伊（Julien d'Huy）利用统计学与系统发生学工具，重建了该故事的旧石器时代版本：“有一种动物，是一种有角的食草动物，尤其是麋鹿。一名人类追逐这只偶蹄动物。狩猎通向天空或发生在天空中。当该动物被转化为星座时，它仍然是活着的。它形成了北斗七星。”$^{[35]}$
\subsubsection{阿拉伯民间传统}  
尽管伊斯兰教出现之前的阿拉伯人将更大的星座——大熊座——识别为一只熊，这或许受到了希腊文化的影响，但他们在传统上始终将北斗七星与小熊座视为一对对应之物。二者都被想象为送葬队伍，其中“斗”的部分被视为灵柩，而“柄”则是随行的哀悼者队列。北斗七星被称为 banāt an-naʿsh al-kubrā，字面意为“灵柩的年长女儿们”。然而这里的“女儿”指的是隶属于灵柩之物，即哀悼者，因此更恰当的译法是“较大的送葬队伍”，而小熊座则被称为“较小的送葬队伍”。另有一种传说认为，灵柩上躺着的是队伍中众人的父亲，他名叫 Naʿash，被 Al-judayy（北极星的阿拉伯名称）所杀，而这支送葬队伍正是在追逐凶手。$^{[36]}$
\subsubsection{希腊—罗马传统}  
在希腊神话中，宙斯（诸神之王，在罗马神话中称为朱庇特）爱慕一位名叫卡利斯托（Callisto）的年轻女子，她是阿耳忒弥斯的宁芙（罗马人称其为狄安娜）。宙斯嫉妒的妻子赫拉（罗马神话中的朱诺）发现卡利斯托因被宙斯强暴而生下了儿子阿耳卡斯（Arcas），于是将卡利斯托变成一只熊作为惩罚。$^{[37]}$处于熊形态的卡利斯托后来与她的儿子阿耳卡斯相遇，阿耳卡斯几乎将这只熊刺杀。为了避免悲剧发生，宙斯将二人一同送上天空，使卡利斯托化为大熊座，而阿耳卡斯化为牧夫座。奥维德称大熊座为“帕拉西亚之熊”（Parrhasian Bear），因为卡利斯托来自阿卡迪亚的帕拉西亚地区，该神话正发生于此。$^{[38]}$

希腊诗人阿拉托斯将该星座称为 Helikē（意为“旋转”或“缠绕”），因为它围绕天极旋转。《奥德赛》中指出，它是唯一一个永不沉入地平线之下、并且“沐浴在大洋波涛之中”的星座，因此常被用作导航的天球参照物。$^{[39]}$它还被称为“车”（Wain）或“Plaustrum”，这是一个拉丁词，指由马牵引的货车。$^{[40]}$
\subsubsection{印度教传统}  
在印度教中，对大熊座／北斗七星／巨熊座的最早记载被称为“七仙人”（Saptarshi），其中每一颗恒星分别代表一位七仙人（Rishi），即：婆利古（Bhrigu）、阿特里（Atri）、安吉拉斯（Angiras）、婆悉吒（Vasishtha）、普拉斯提亚（Pulastya）、普拉哈（Pulaha）与克拉图（Kratu）。这一记载见于《梨俱吠陀》（约公元前1500—1200年），该文献是人类已知最古老的文本之一。

古印度文献中的相关记载包括：
\begin{enumerate}
\item 《梨俱吠陀》（第一曼荼罗，第24赞歌，第10节），在宇宙秩序与天体意义的语境中提及七仙人。
\item 《摩诃婆罗多》（约公元前400年至公元400年），论及“七仙人星群”（Saptarishi Mandal）作为航行与指引的恒星。
\item 《往世书》（如《毗湿奴往世书》《梵天卵往世书》等），将七仙人描述为掌握宇宙智慧的神圣贤者。
\end{enumerate}
至于小熊座，其并未在早期吠陀文献中被明确提及，但在后来的天文学著作中得到了确认，例如：
\begin{enumerate}
\item 《吠檀伽天文书》（Vedanga Jyotisha，约公元前1400—1200年）
\item 《苏利耶悉檀多》（Surya Siddhanta，约公元4世纪）
\end{enumerate}
关于大熊座前方两颗恒星指向北极星的事实，被解释为毗湿奴赐予少年圣者德鲁瓦（Dhruva）的恩典。$^{[41]}$

因此，《梨俱吠陀》构成了关于大熊座的最早文字记载，而小熊座则是在后期天文学传统中逐渐获得重要地位。
\subsubsection{在犹太教与基督教中}  
大熊座可能在《圣经·约伯记》中被提及，该书成书时间一般认为介于公元前7世纪至公元前4世纪之间，但这一对应关系常存在争议。$^{[42]}$
\subsubsection{东亚传统}  
在中国与日本，北斗七星被称为“北斗”（中文：běidǒu，日文：hokuto）。在古代，七颗恒星各自都有专名，多源自中国古代星名体系：

\begin{itemize}
\item “枢”（樞，中文：shū；日文：sū）对应天枢（Dubhe，大熊座$\alpha$星）  
\item “璇”（璇，中文：xuán；日文：sen）对应天璇（Merak，大熊座$\beta$星）  
\item “玑”（璣，中文：jī；日文：ki）对应天玑（Phecda，大熊座$\gamma$星）  
\item “权”$^{[43]}$（權，中文：quán；日文：ken）对应天权（Megrez，大熊座 $\delta$ 星）  
\item “玉衡”（中文：yùhéng；日文：gyokkō）对应玉衡（Alioth，大熊座$\epsilon$ 星）  
\item “开阳”（開陽，中文：kāiyáng；日文：kaiyō）对应开阳（Mizar，大熊座$\zeta$ 星）  
\item 摇光（Alkaid，大熊座$\eta$星）则拥有多个别名：“剑”（劍，中文：jiàn；日文：ken，为“剑先”劍先的简化形式）、“摇光”（搖光，中文：yáoguāng；日文：yōkō），以及“破军星”（破軍星，中文：pójūn xīng；日文：hagun sei），因为在传统观念中，军队若向该星方向行军被视为不祥之兆。$^{[44]}$
\end{itemize}
在神道教中，大熊座中最亮的七颗恒星隶属于天之御中主神（Ame-no-Minakanushi），被视为最古老且最具力量的神祇。

在韩国，该星座被称为“北方七星”。相关神话讲述了一位寡妇与其七个儿子的故事：寡妇爱上一位鳏夫，但通往其住所需横渡一条河流。七个儿子体恤母亲，在河中放置踏石。母亲并不知情，却为这些踏石祝福；待七个儿子去世后，他们化为天上的星座。
\subsubsection{美洲原住民传统}  
易洛魁人将玉衡（Alioth）、开阳（Mizar）与摇光（Alkaid）视为三名猎人，正在追逐巨熊。按照其中一种神话版本，第一位猎人（玉衡）手持弓箭，准备射杀巨熊；第二位猎人（开阳）肩扛一口大锅——即其伴星辅星（Alcor）——用于烹煮猎物；第三位猎人（摇光）则拖拽着一捆柴火，用以在锅下生火。

拉科塔人称该星座为 Wičhákhiyuhapi，意为“巨熊”。

根据托马斯·莫顿在《新英格兰迦南》中的记载，万帕诺亚格人（阿尔冈昆语族）将大熊座称为 “maske”，意为“熊”。$^{[45]}$

瓦斯科—威什拉姆族原住民则将该星座解释为五只狼与两只熊，它们被郊狼神留在了天空中。$^{[46]}$
\subsubsection{日耳曼传统}  
在北欧异教传统中，北斗七星被称为 Óðins vagn，即“奥丁之车”。同样地，奥丁在诗性称谓（Kennings）中也被称为 vagna verr（“车之守护者”）或 vagna rúni（“车之密友”）。$^{[47]}$
\subsubsection{乌拉尔语系传统}  
在芬兰语中，这一星群有时以其古老的芬兰名称 Otava 被称呼。该名称在现代芬兰语中的含义几乎已经被遗忘，其原意是“鲑鱼拦网”。古代芬兰人认为，熊（Ursus arctos）是被装在金篮中从大熊座降到地上的；当一头熊被猎杀后，人们会将熊头放置在树上，以便让熊的灵魂返回大熊座。

在北欧的萨米语中，这一星座的一部分（即北斗七星去除天枢与天璇）被视为伟大猎人 Fávdna 的弓（猎人对应的恒星为大角星 Arcturus）。在主要的萨米语言——北萨米语中，它被称为 Fávdnadávgi（“Fávdna 之弓”），或简称为 dávggát（“弓”）。这一星座在萨米族国歌中占据重要地位，歌词开头为 Guhkkin davvin dávggaid vuolde sabmá suolggai Sámieanan，其意为“在遥远的北方，在弓之下，萨米之地缓缓显现”。“弓”在萨米族关于夜空的传统叙事中具有核心意义，其中多位猎人试图追逐 Sarva——巨大的驯鹿星座，它几乎占据了半个天空。根据传说，Fávdna 每晚都准备拉弓射击，但他迟疑不决，因为他可能会击中被称为 Boahji（“铆钉”）的北极星（Stella Polaris），这将导致天空崩塌并终结世界。$^{[48]}$
\subsubsection{东南亚传统}  
在缅甸语中，Pucwan Tārā（ပုဇွန် တာရာ，[bəzʊ̀ɴ tàjà]）是一个由大熊座头部与前肢恒星组成的星座名称；pucwan（ပုဇွန်）是对甲壳类动物的统称，如对虾、虾、蟹、龙虾等。

在爪哇语中，该星座被称为 “lintang jong”，意为“舢板星座”；在马来语中则被称为 “bintang jong”。$^{[49]}$
\subsubsection{神秘学传统}  
在神智学（Theosophy）中，人们相信昴宿星团的七颗星将来自银河洛格斯（Galactic Logos）的七道灵性能量汇聚并传导至大熊座的七颗恒星，再传至天狼星，随后传至太阳，再到地球之神（Sanat Kumara），最终通过七道光线的七位大师传递给人类。$^{[50]}$
\subsection{在文化中}  
大熊座曾被荷马、斯宾塞、莎士比亚、丁尼生等诗人提及，也出现在费德里科·加西亚·洛尔迦的《献给月亮之歌》中。$^{[51]}$古代芬兰诗歌同样提到这一星座，它还出现在文森特·梵高的绘画作品《罗讷河上的星夜》中。$^{[52][53]}$
\subsection{图形可视化}
在欧洲的星图中，这一星座通常被描绘为：北斗七星中构成“斗”的四颗星形成熊的身体，而由其余三星构成的“柄”则被画成一条修长的尾巴。然而，熊并没有长尾巴，因此犹太天文学家将五帝座一（Alioth）、开阳（Mizar）和摇光（Alkaid）视为跟随母熊的三只幼熊；而在美洲原住民的观念中，这三颗星则被看作是三名猎人。
\begin{figure}[ht]
\centering
\includegraphics[width=6cm]{./figures/4af729d04d057a7a.png}
\caption{H.~A.~Rey 为大熊座提出的另一种星群构想，可以认为赋予了其类似北极熊那样更长的头部与颈部，如这张从左侧拍摄的照片所示。} \label{fig_UrMa_6}
\end{figure}
著名儿童读物作家 H.~A.~Rey 在其 1952 年出版的著作《The Stars: A New Way to See Them》（ISBN 0-395-24830-2）中，对大熊座提出了一种不同的星群想象方式。在他的构想中，“熊”的形象被重新定向：以摇光（Alkaid）作为熊鼻子的尖端，而北斗七星的“柄”部分则构成了熊头顶部与颈部的轮廓，并向后延伸至肩部，从而使其整体形态更接近具有较长头部和颈部的北极熊。
\begin{figure}[ht]
\centering
\includegraphics[width=14.25cm]{./figures/8c59e3b77bd1125f.png}
\caption{} \label{fig_UrMa_7}
\end{figure}
大熊座还被描绘为“星犁”（Starry Plough）：这是爱尔兰工党旗帜的象征，由詹姆斯·康诺利于1916年为爱尔兰公民军所采用，旗面为蓝底，其上绘有该星座；它也出现在美国阿拉斯加州州旗上，以及瑞典贝尔纳多特王朝纹章的一种变体中。西班牙马德里自治区旗帜上红底的七颗星，可能象征“犁”这一星群（或小熊座）。同样的解释也适用于西班牙首都马德里市徽章中蓝色饰边（bordure azure）上描绘的七颗星。
\subsection{另见}
\begin{itemize}
\item 大熊座（中国天文学）
\item 小熊座
\item 南十字座
\item 天球制图学
\item 星座家族
\item 已废弃的星座
\item 按星座列出的恒星列表
\item 约翰内斯·赫维留斯所列星座
\item 拉卡伊所列星座
\item 彼得鲁斯·普朗修斯所列星座
\item 托勒密所列星座
\end{itemize}
\subsection{注释}

a.托勒密以希腊语将该星座命名为 Ἄρκτος μεγάλη（Arktos Megale，意为“大熊”）。小熊座则被称为 Ἄρκτος μικρά（Arktos Mikra）。\\
\subsection{参考文献}
 柯克帕特里克，J·戴维；马罗科，费德里科等（2024年4月）。“基于全天空20秒差距范围内约3600颗恒星与棕矮星普查的初始质量函数”. 天体物理学杂志增刊系列. 271 (2): 55. arXiv:2312.03639. 文献标识码:2024ApJS..271...55K. doi:10.3847/1538-4365/ad24e2.
 “钱德拉：大熊座”. chandra.harvard.edu. 检索日期2022-07-07.
 “星座 | 宇宙”. astronomy.swin.edu.au. 检索日期2022-07-07.
 天蝎座–狐狸座.
 “星座”。国际天文学联合会。2013年6月4日从原文存档。检索日期：2019年8月16日。
 读者文摘协会（2005年8月）。地球与宇宙。读者文摘协会有限公司。ISBN978-0-276-42715-2。于2021年4月15日存档自原始网页。检索日期：2016年11月7日
 “查尔斯的酒馆”。归档于2017年12月1日自原始页面。检索于2017年11月23日
 安德烈·G·博尔德洛（2013年10月22日）。夜空的旗帜：当天文学邂逅民族自豪感。施普林格科学与商业媒体。第131页及以后。ISBN978-1-4614-0929-8。于2021年4月15日存档自原始网页。检索于2019年6月7日
 詹姆斯·B·凯勒（2011年7月28日）。恒星及其光谱：光谱序列导论。剑桥大学出版社。第241页起。ISBN978-0-521-89954-3。于2021年4月15日存档自原始网页。检索于2019年6月7日
 吉姆·凯勒（2009-09-16），“恒星：‘阿利奥特’”“。”于2019-12-11从原始网页存档。检索日期：2019-06-07
 马克·R·查特兰德（1982）。《天导：业余天文学家实地指南》。黄金出版社。文献标识码:1982sfga.book.....C. ISBN 978-0-307-13667-1. 已存档于2021年4月15日自原始网页存档。已检索2019-06-07.
 里德帕斯，第136页。
 伊恩·里德帕斯。“星图——大熊座”.
 “特里维克拉马”.全天空百科全书国际天文学联合会恒星命名工作组。检索日期：2025年10月15日
 莱维 2005, 第67页。
 “行星47乌玛b”.系外行星百科全书。2012年7月11日。于2023年11月11日存档自原始网页。检索日期：2012年7月15日
 “行星47乌玛c”。系外行星百科全书。2012年7月11日。于2023年11月11日存档自原始网页。检索日期：2012年7月15日
 “行星47乌玛 d”.系外行星百科全书。2012年7月11日。于2023年11月11日从原始网页存档。检索日期：2012年7月15日
 “特拉华州事实与符号——特拉华州其他符号”.delaware.gov已于2019年3月28日存档自原始页面。检索日期：2019年6月7日
 威尔金斯，杰米；邓恩，罗伯特（2006）。300个天体：宇宙视觉参考（第1版）。纽约布法罗：萤火虫图书。ISBN 978-1-55407-175-3.
 曹，Y；卡斯利瓦尔，M. M；麦凯，A；布拉德利，A（2014）。“对M82中超新星的分类：一颗年轻的、呈红色的Ia型超新星”。天文学家电报. 5786: 1. 文献标识码:2014ATel.5786....1C.
 莱维 2005, 第129–130页。
 塞罗尼克，加里（2012年7月）。“M101：一只巨型星系”。天空与望远镜. 124 (1): 45. Bibcode:2012S&T...124a..45S.
 哈伊杜科娃，M.；内斯鲁尚，L.（2020）。“仙女座χ流星雨与1月大熊座α流星雨：与C/1992 W1（大下田）彗星相关的新流星雨及可能的流星雨”。国际行星科学杂志. 351 113960. 文献标识码:2020Icar..35113960H. DOI:10.1016/j.icarus.2020.113960. S2CID 224889918.
 哈伊杜科娃，玛丽亚；内斯鲁尚，卢博什（2021-06-28）。未知的彗星C/1992 W1（大下）和C/1853 G1（施韦策）伴生流星雨. 欧洲行星科学大会2021. 文献标识码:2021EPSC...15..214H. DOI:10.5194/epsc2021-214.
 詹尼斯肯斯，彼得（2012年9月）：“绘制流星体轨道：新流星雨被发现”。天空与望远镜: 23.
 加伊多什，斯特凡（2008-06-01）。“寻找大熊座十月流星雨的过去迹象”。地球、月球与行星. 102 (1): 117–123. 文献标识码:2008EM&P..102..117G. DOI:10.1007/s11038-007-9196-9. ISSN 1573-0794. S2CID 123555821.
 “十月大熊座流星雨——仰望星空”. blogs.nasa.gov. 2011年10月18日. 检索日期2023-03-30.[permanent dead link]
 劳克林，格雷格（2013年5月）。“世界为何会失衡”。天空与望远镜. 125 (5): 29. Bibcode:2013S&T...125e..26L.
 马洛里，J.P.；亚当斯，D.Q.（2006年8月）。“第8.5章：原始印欧人的自然地理环境”。牛津原始印欧语与原始印欧世界导论。英国牛津：牛津大学出版社。第131页。ISBN 9780199287918. OCLC 139999117. 在印欧语系中，重建得最稳固的星座当属大熊座。在希腊语和梵语中，它被称为“熊”（第9章；拉丁语此处可能是借自其他语言），不过即便是后一种称谓也受到了质疑。
 里德帕斯，伊恩.“托勒密的天球图：首次印刷版，1515年”。检索日期：2022年11月15日
 艾伦，R. H.（1963）。恒星名称：其传说与含义（重印版）。纽约，纽约：多佛出版公司，第207–208页。ISBN978-0-486-21079-7。检索日期：2010-12-12：CS1维护：忽略ISBN错误（链接)
 吉本，威廉·B.（1964）。“北美星象传说中的亚洲相似之处：大熊座”。美国民俗学杂志. 77 (305): 236–250. DOI:10.2307/537746. JSTOR 537746.
 布拉德利·E·谢弗希腊星座的起源：大熊座是否在猎人游牧民族于一万三千多年前首次抵达美洲之前就已命名？, 科学美国人, 2006年11月，评述于《希腊星座的起源》 存档2017-04-01 于互联网档案馆； 尤里·别列兹金,宇宙狩猎：西伯利亚—北美神话的变体 存档2015-05-04 于互联网档案馆. 民俗学, 31, 2005: 79–100.
 杜伊·朱利安《星空中的熊：一项史前神话的系统发生学研究》 已归档2021-12-20于互联网档案馆, 西南史前, 20 (1), 2012: 91–106; 柏柏尔天空中的宇宙狩猎：旧石器时代神话的系统发生学重构 已归档2020-05-28于互联网档案馆, AARS期刊, 15, 2012.
 “莱恩词典 阿拉伯语到英语 第2816页”.
 “大熊座，大熊”. 伊恩·里德帕斯的星图故事”.
 奥维德，《英雄书信集》（格兰特·肖弗曼译）第18封信
 荷马，《奥德赛》，第五卷，273
 “阿皮安努斯对大熊座和小熊座的描绘”. 伊恩·里德帕斯的星图故事.
 马哈德夫·哈里拜·德赛（1973）。与甘地的日常点滴：秘书日记。萨尔瓦·塞瓦·桑格出版社。已存档自原始网页于2022年5月13日。检索日期：2021-01-06.
 博特韦克，G·约翰内斯，编（1994）。旧约神学辞典，第7卷。Wm. B. 埃尔德曼斯出版社。第79–80页。ISBN978-0-8028-2331-1。于2022年4月7日存档自原版。检索日期：2019年7月30日
 《中国星区、星群与恒星名称英汉对照表》。香港太空馆于2018年12月17日存档自原始页面。检索日期：2018年12月17日
 《伴仙集》，由忍者大师藤林安武于1676年撰写，书中多次提及这些星辰，并在第8册第17卷中展示了一幅传统的北斗七星图，该卷内容涉及天文学与气象学（摘自阿克塞尔·马祖埃的译本）。
 托马斯，莫顿（1883年9月13日）。托马斯·莫顿的全新英格兰迦南。由王子学会出版。OL 7142058M.
 克拉克，埃拉·伊丽莎白（1963）。太平洋西北部的印第安传说。加州大学出版社。于2022-05-13从原始版本存档。检索日期：2019-05-01。
 克利斯比，理查德; 维格弗松，古德布兰德 (1874). 冰岛语-英语词典。牛津：克拉伦登出版社。第674页。
 自然科学中心网：星空（https://www.naturfagsenteret.no/c1515376/binfil/download2.php?tid=1509706)
 伯内尔，A.C.（2018）。霍布森-乔布森：英印俗语及短语词典。劳特利奇出版社。第472页。ISBN 9781136603310.
 Baker, Dr. Douglas The Seven Rays:Key to the Mysteries 1952
 "Canción para la luna - Federico García Lorca - Ciudad Seva". Archived from the original on 2015-05-10. Retrieved 2015-08-16.
 青蛙（2018-01-30）。“神话”. 人文. 7 (1): 14. doi:10.3390/h7010014. ISSN 2076-0787.
 克莱森，霍利斯（2002）。“展览评论：‘有些事物会结出果实吗？’——见证梵高与高更之间的纽带”。艺术公报. 84 (4): 670–684. DOI:10.2307/3177290. ISSN 0004-3079. JSTOR 3177290.
 “H.A.雷伊为大熊座所作星群的存档版”。已于2014年4月7日从原始页面存档。
\subsubsection{参考书目}
列维，大卫·H.（2005）。深空天体。普罗米修斯图书。ISBN 978-1-59102-361-6.
汤普森，罗伯特；汤普森，芭芭拉（2007）。天文奇观图解指南：从新手到资深观测者。奥莱利媒体公司。ISBN 978-0-596-52685-6.
延伸阅读
伊恩·里德帕斯和威尔·蒂里昂（2007）。恒星与行星指南, 科林斯，伦敦。ISBN 978-0-00-725120-9。普林斯顿大学出版社，普林斯顿。ISBN 978-0-691-13556-4.
外部链接
